%%% Results - EBNA1 case study

To illustrate IMPG's capacity for rapid population-scale exploration, we investigated a tandem repeat cluster at chromosome 11q23 that Li et al.\ \cite{Li2023} identified as bound by the Epstein--Barr virus EBNA1 protein, inducing chromosomal fragile site breakage. The original study examined four long-read genomes---two Ashkenazi and two Chinese---and reported a $\sim$2.6-fold difference in repeat copy number between the two groups. Using IMPG, we queried this $\sim$21\,kb locus across all 466 HPRCv2 haplotypes in seconds (Figure~\ref{fig:ebv}). Across five superpopulations, repeat region lengths range from 19,941 to 24,468\,bp (1.2-fold) and EBNA1-like motif counts from 446 to 551 per haplotype (1.2-fold), with identical median values across all superpopulations. The apparent $\sim$2.6-fold population difference reported from four genomes is not supported at population scale: variation at this locus is continuous across ancestries, not a discrete two-group difference, illustrating how small sample sizes can mischaracterize structural variation at clinically relevant loci.
